\documentclass[11pt]{article}
\usepackage{fontspec}
\usepackage[french]{babel}
\usepackage{amsmath}
\usepackage{amssymb}
\usepackage{amsthm}
\usepackage{mathrsfs}
\usepackage{tikz-cd}
\usepackage{enumitem}

\newcommand{\printedpage}[1]{\marginpar{\footnotesize p.~#1}}
\newcommand{\msp}[1]{\marginpar{\footnotesize [#1]}}
\newcommand{\circled}[1]{\tikz[baseline=(c.base)]\node[draw,circle,inner sep=1pt](c){#1};}
\newcommand{\deff}{\overset{\text{déf}}{=}}

\begin{document}

\printedpage{140}
% head: A. Grothendieck : Dérivateurs, Ch. XII, Édition des Universités de Montpellier II et Paris VII

\section*{5 Axiomatique des $W$-types d'homotopie relatifs non localement constants $\mathrm{Ho}^{\mathcal{L}}_W(S)$.}

On n'a rien dit jusqu'à ce point d'éventuels $W$-types d'homotopie relatifs (sur un objet $S$ de $\mathcal{M}$) qui ne seraient \textit{pas} ``localement constants''. Il en était question pourtant dans loc. cit. (VI, théorème-scholie 2, partie II du théorème pages 109-113.) Je profite de l'occasion pour revenir là dessus. Il faut introduire, dans chaque $\mathcal{M}/S$, un deuxième localiseur
\[
W^0_S \subseteq \mathrm{Fl}(\mathcal{M}/S), \qquad W^0_S \subseteq W_S
\]
(plus fin que $W_S$), ce qui permettra de définir
\msp{page 219}
\[
\mathrm{Ho}_W(S) \deff (\mathcal{M}/S)(W^0_S)^{-1}.
\]
Qu'il n'y ait pas de choix \textit{canonique} qui s'impose ici, déduit de la seule structure $(\mathcal{M},W)$, est claire sur l'exemple de $(\mathrm{Cat},W)$, auquel cas on trouve \textit{deux} choix symétriques, transformés l'un en l'autre par l'involution
\[
X \longmapsto X^\circ : (\mathrm{Cat},W) \xrightarrow{\sim} (\mathrm{Cat},W),
\]
à savoir $W^{\mathrm{g}}_S$ et $W^{\mathrm{d}}_S$. Le choix de l'un ou de l'autre, disons (pour fixer les idées) de $W^{\mathrm{g}}_S$, est associé à une notion d'\textit{isomorphisme local} (ou de `coisomorphisme local' dans le cas de $W^{\mathrm{d}}_S$) que je voudrais axiomatiser. Il semble bien que pour les principales catégories de modèles pour Hot (et pour ses variantes), on dispose de façon assez naturelle d'une telle notion d'isomorphisme local (sans l'embarras, le plus souvent, d'un automorphisme involutif fournissant une notion symétrique \dots).

On supposera donc, comme dans loc. cit. p. 94-96, donné une classe
\[
\mathcal{L} \subseteq \mathrm{Fl}(\mathcal{M})
\]
de morphismes, jounant le rôle de `flèches de % [sic] jounant -> jouant, faute de frappe
\msp{page 220}
localisation' ou `d'isomorphismes locaux', satisfaisant les conditions suivantes.

\begin{enumerate}
\item[$(*)$] $\mathcal{L}$ contient les isomorphismes, est stable par composition, et par changement de base.
\end{enumerate}

On pose, pour $S$ dans $\mathcal{M}$,
\[
\mathcal{L}_S \subseteq \mathcal{M}/S \qquad \text{sous-catégorie pleine des } S'/S \text{ telle que } S' \longrightarrow S \text{ \texttt{[soit]} dans } \mathcal{L}.
\]
On définit alors pour tout $S$ dans $\mathcal{M}$
\[
\begin{aligned}
&W^{\mathcal{L}}_S \subseteq \mathrm{Fl}(\mathcal{M}/S)\\
&(X \xrightarrow{u} Y) \in W^{\mathcal{L}}_S \iff \forall\, S' \in \mathcal{L}_S,\ (u_{S'} : X \times_S S' \longrightarrow Y \times_S S') \text{ est dans } W.
\end{aligned}
\]
% foot: Version 30 mars 2003

\printedpage{141}
% head: A. Grothendieck : Dérivateurs, Ch. XII, Édition des Universités de Montpellier II et Paris VII

C'est donc une variante du localiseur $W^{\mathrm{u}}_S$, où maintenant on ne prend comme changements de base que les flèches `$\mathcal{L}$-localisantes', i.e. $\in \mathcal{L}$. On a donc
\[
W^{\mathrm{u}}_S \subseteq W^{\mathcal{L}}_S \subseteq W_S
\]
en tant que parties de $\mathrm{Fl}(\mathcal{M}/S)$, et de même dans toute sous-catégorie pleine $\mathcal{F}$ de $\mathcal{M}/S$. Dans le cas de $\underline{\mathrm{Fib}}_W S$, moyennant l'axiome \circled{b} ci-dessus ($W^{\mathrm{u}}_S = W_S$ dans $\underline{\mathrm{Fib}}_W S$), les trois localiseurs coïncident.

On pose
\[
\mathrm{Ho}^{\mathcal{L}}_W(S) \text{ (ou } \mathrm{Ho}_W(S)\text{)} \ \deff (\mathcal{M}/S)(W^{\mathcal{L}}_S)^{-1},
\]
et on a donné dans loc. cit. (partie II du théorème-scholie) un résultat intéressant sur cet invariant, pour $S$ fixé. Mon propos
\msp{page 221}
présent est de voir ce qui se passe pour $S$ variable. Visiblement, une notion de $W$-type d'homotopie relatif sur un objet quelconque $S$ de $\mathcal{M}$ n'est satisfaisante, que si elle permet des images inverses par des flèches \textit{quelconques}
\[
f : S' \longrightarrow S
\]
dans $\mathcal{M}$, qui doivent définir absolument
\[
f^{*\mathcal{L}}_W \text{ ou } f^*_W : \mathrm{Ho}^{\mathcal{L}}_W(S) \longrightarrow \mathrm{Ho}^{\mathcal{L}}_W(S').
\]
On est donc confronté à la même perplexité que tantôt - et ce n'est \textit{pas,} cette fois, l'introduction de $\underline{\mathrm{Fib}}_W(S)$ qui va sauver la situation, car
\[
(\underline{\mathrm{Fib}}_W S)(\underbrace{W^{\mathcal{L}}_S}_{\substack{= W_S \\ = W^{\mathrm{u}}_S}})^{-1} \longrightarrow (\mathcal{M}/S)(W^{\mathcal{L}}_S)^{-1}
\]
n'a aucune envie, cette fois, d'être une équivalence de catégories, mais bien plutôt une inclusion (pleinement fidèle) des $W$-types d'homotopie relatifs constants sur $S$, dans les $W$-types d'homotopie quelconques. Il faut visiblement un axiome qui remplace l'axiome `de la fibre homotopique'.

\msp{page 222}

Dans le cas de Cat, la classe $\mathcal{L}$ des isomorphismes locaux (disons) permet de définir, comme ci-dessus, le localiseur $W^{\mathcal{L}}_S = W^{\mathrm{g}}_S$ dans les $\mathrm{Cat}/S$, comme formé des $u : X \longrightarrow Y$ sur $S$ qui sont dans $W$ et le restent par tout changement de base $S' \longrightarrow S$ qui est dans $\mathcal{L}$. Mais il y a lieu de \textit{saturer} $\mathcal{L}$, en regardant la classe $\mathcal{L}^\wedge$ des flèches $f : S \longrightarrow S'$ de Cat telles que le changement de base
\[
f^* : \mathrm{Cat}/S \longrightarrow \mathrm{Cat}/S'
\]
envoie $W^{\mathrm{g}}_S$ dans $W^{\mathrm{g}}_{S'}$. Ce $\mathcal{L}^\wedge$ est tautologiquement stable par composition, et contient les isomorphismes (mieux, contient $\mathcal{L}$ = ensemble des isomorphismes locaux). Il n'est pas
% foot: Version 30 mars 2003
\printedpage{142}
% head: A. Grothendieck : Dérivateurs, Ch. XII, Édition des Universités de Montpellier II et Paris VII
clair, par la définition donnée ici, qu'il soit stable par changement de base - on va renforcer la définition en conséquence, en exigeant, pour que $f : S' \longrightarrow S$ soit dans $\mathcal{L}^\wedge$, que pour tout $T$ sur $S$, considérons $f_T : T' = S' \times_S T \longrightarrow T$, $f^*_T : \mathrm{Cat}/T \longrightarrow \mathrm{Cat}/T'$ soit compatible avec les localiseurs $W^{\mathrm{g}}_T$, $W^{\mathrm{g}}_{T'}$. Alors $\mathcal{L}^\wedge$ est (tout comme $\mathcal{L}$) stable par composition et par changement de base, et de plus il contient $\mathcal{L}$. De plus, pour tout $S$, on a $W^{\mathcal{L}}_S = W^{\mathcal{L}^\wedge}_S$. Ce petit sorite pourrait se faire bien sûr pour tout modéliseur $(\mathcal{M},W)$. Mais la chose intéressante ici, c'est que $\mathcal{L}^\wedge$ contient l'ensemble $\mathrm{Liss}_W$ des flèches $W$-lisses. (J'ignore s'il est égal.) On aurait donc pu
\msp{page 223}
aussi bien partir de l'ensemble de flèches $\mathrm{Liss}_W$, pour définir les localiseurs $W^{\mathrm{g}}_S$ comme
\[
W^{\mathrm{g}}_S = W^{\mathrm{Liss}_W}_S \qquad \text{au lieu de } W^{\mathrm{g}}_S = W^{\mathrm{Isoloc}}_S,
\]
ce qui est équivalent.

D'autre part, en travaillant avec les localisés $\mathrm{HOT}^\circ_W(S) = (\mathrm{Cat}/S)(W^{\mathrm{g}}_S)^{-1}$, le rôle de $\underline{\mathrm{Fib}}_W S$ pour la théorie des $\mathrm{HOT}^{\mathrm{lc}}_W(S)$ est joué ici par $\underline{\mathrm{Prop}}_W(S)$. On voit donc que dans la théorie des $W$-types d'homotopie relatifs non localement constants, mode $\mathrm{HOT}^\circ_W$, les ensembles de flèches $\mathrm{Liss}_W$ et $\mathrm{Prop}_W$ jouent un rôle en quelque sorte complémentaire : les flèches $f : S' \longrightarrow S$ de $\mathrm{Liss}_W$ sont (pratiquement) celles qui permettent directement de définir
\[
f^{*\circ}_W : \mathrm{HOT}^\circ_W(S) \longrightarrow \mathrm{HOT}^\circ_W(S')
\]
comme induit par
\[
f^* : (\mathrm{Cat}/S, W^{\mathrm{g}}_S) \longrightarrow (\mathrm{Cat}/S', W^{\mathrm{g}}_{S'}),
\]
tandis que les flèches de $\mathrm{Liss}_W$ permettent de définir pour tout $S$ dans Cat la sous-catégorie pleine
\[
\underline{\mathrm{Prop}}_W S \subseteq \mathrm{Cat}/S,
\]
donnant lieu à une équivalence de catégories
\[
\left\{
\begin{aligned}
&(\underline{\mathrm{Prop}}_W S)(W^{\mathrm{u}}_S)^{-1} \xrightarrow{\sim} \underbrace{(\mathrm{Cat}/S)(W^{\mathrm{g}}_S)^{-1}}_{= \mathrm{HOT}^0_W(S)} % ?? uncertain: exposant de HOT sous l'accolade -- 0 (zéro) ou \circ
\\
&\text{avec}\\
&W^{\mathrm{u}}_S = W^{\mathrm{g}}_S \text{ sur } \underline{\mathrm{Prop}}_W S
\end{aligned}
\right.
\]
(ce qui visiblement est le substitut annoncé de ``l'axiome de la fibre $W$-homotopique''). Cette
\msp{page 224}
équivalence, qui permet de décrire $\mathrm{Hot}^\circ_W(S)$ en termes de $\underline{\mathrm{Prop}}_W S \subseteq \mathrm{Cat}/S$ et d'un localiseur $W^{\mathrm{u}}_S$ \textit{stable par} \textbf{tout} \textit{changement de base,} permet d'autre part de définir la loi contravariante attendue pour $\mathrm{HOT}^\circ_W(S)$.
% foot: Version 30 mars 2003

\printedpage{143}
% head: A. Grothendieck : Dérivateurs, Ch. XII, Édition des Universités de Montpellier II et Paris VII

D'autre part, quand on échange le rôle de $\mathrm{Liss}_W$ et $\mathrm{Prop}_W$, on trouve la théorie complémentaire du $\mathrm{HOT}_W(S)$,
\[
\left\{
\begin{aligned}
&\mathrm{HOT}_W(S) \deff (\mathrm{Cat}/S)(W^{\mathrm{d}}_S)^{-1} \xleftarrow{\sim} (\underline{\mathrm{Liss}}_W S)(W^{\mathrm{u}}_S)^{-1},\\
&\text{sur } \underline{\mathrm{Liss}}_W S,\ W^{\mathrm{d}}_S = W^{\mathrm{u}}_S,
\end{aligned}
\right.
\]
la description de $\mathrm{HOT}_W(S)$ en termes de $\underline{\mathrm{Liss}}_W S$ et de $W^{\mathrm{u}}_S$ permet d'en définir le caractère contravariant en $S$, pour \textit{toute} flèche $f : S' \longrightarrow S$ dans Cat. Tandis que si on veut pouvoir passer aux catégories de fractions sur $\mathrm{Cat}/S$ et $\mathrm{Cat}/S'$ tout entier (via $W^{\mathrm{d}}_S$, $W^{\mathrm{d}}_{S'}$), il faut se limiter aux changements de base $W$-propres.

Du coup, me vient la perplexité si cette complémentaire, voire cette dualité, entre deux ensembles de flèches, les `lisses' et les `propres', ne serait pas un fait général qui se retrouvait tout au moins dans les catégories de modèles proche de la description de HOT ou d'un $\mathrm{Hot}_W$.

\msp{page 225}

Il se pourrait bien que dans $\boldsymbol{\Delta}^\wedge$ on ait une dualité parfaite entre ces deux espèces de notions (il faudra que j'examine la chose), peut-être même dans toute catégorie $\mathcal{A}^\wedge$ (mais dans $\boldsymbol{\Delta}^\wedge$ il y a un automorphisme involutif $K \longmapsto K^\circ$, correspondant au $X \longmapsto X^\circ$ dans Cat \dots). Et après tout, même dans le contexte des espaces topologiques, les notions de propreté d'une part, d'asphéricité locale d'autre part, pour une flèche $X \longrightarrow Y$, sont bien définies, et pourraient bien s'ajuster comme le font les notions correspondantes dans Cat, et (qui sait) donner lieu à deux versions complémentaires, ou `duales' l'une de l'autre, d'une notion de type d'homotopie relatif.

N'en ayant pas le cœur net, je suis un peu indécis comment présenter une description axiomatique de la théorie des $\mathrm{Ho}^{\mathcal{L}}_W(S)$, pour une catégorie de modèles générale $(\mathcal{M},W)$. Dois-je me donner d'emblée
\[
\mathrm{Liss}, \mathrm{Prop} \subseteq \mathrm{Fl}(\mathcal{M}),
\]
stables par composition, par changement de base, contenant les isomorphismes, et tels que (peut-être)
\[
\mathrm{Liss} \cap \mathrm{Prop} \subseteq \mathrm{Fib}_W,
\]
\msp{page 226}
mais surtout ils permettent de définir, pour $S$ de $\mathcal{M}$
\[
\left\{
\begin{array}{lcccl}
W^{\mathrm{g}}_S & \deff & W^{\mathrm{Liss}}_S & \subseteq & \mathcal{F}(\mathcal{M}/S)\\
W^{\mathrm{d}}_S & \deff & W^{\mathrm{Prop}}_S & \subseteq & \mathcal{F}(\mathcal{M}/S)
\end{array}
\right.
\]
avec les propriétés suivantes :

\begin{enumerate}
\item[{\circled{a}}] Si $f : S' \longrightarrow S$ est dans Prop (resp. Liss), alors $f^* : \mathcal{M}/S \longrightarrow \mathcal{M}/S'$ transforme $W^{\mathrm{d}}_S$ dans $W^{\mathrm{d}}_{S'}$ (resp. $W^{\mathrm{g}}_S$ dans $W^{\mathrm{g}}_{S'}$),
\end{enumerate}
% foot: Version 30 mars 2003
\printedpage{144}
% head: A. Grothendieck : Dérivateurs, Ch. XII, Édition des Universités de Montpellier II et Paris VII
\noindent ce qui est une pure tautologie, et ne concerne pas encore les relations de Prop avec Liss, mais par contre :

\begin{enumerate}
\item[{\circled{b}}]
\[
\left\{
\begin{aligned}
&(\underline{\mathrm{Liss}}S)(W^{\mathrm{u}}_S)^{-1} \xrightarrow{\sim} (\mathcal{M}/S)(W^{\mathrm{d}}_S)^{-1} \text{ \texttt{[est]} une équivalence}\\
&\text{et } W^{\mathrm{u}}_S = W^{\mathrm{d}}_S \text{ sur } \underline{\mathrm{Liss}}S,
\end{aligned}
\right.
\]
et dualement
\[
\left\{
\begin{aligned}
&(\underline{\mathrm{Prop}}S)(W^{\mathrm{u}}_S)^{-1} \xrightarrow{\sim} (\mathcal{M}/S)(W^{\mathrm{g}}_S)^{-1} \text{ \texttt{[est]} une équivalence}\\
&\text{et } W^{\mathrm{u}}_S = W^{\mathrm{g}}_S \text{ sur } \underline{\mathrm{Prop}}S.
\end{aligned}
\right.
\]
\end{enumerate}

De plus, il y a lieu d'énoncer \circled{c} deux propriétés d'adjonction (duales l'une de l'autre), pour lesquelles je renvoie à VI, théorème-scholie 2, p. 108-113. Il se pose ici la question subsidiaire de simplifier les conditions de validité de la formule d'adjonction, un peu hétéroclites dans loc. cit.

Il se pose d'autre part la question de faire la théorie des lois covariantes $f_!$ et $f_*$ pour $\mathrm{Ho}^{\mathrm{lc}}_W(S)$ et pour $\mathrm{Ho}^{\mathrm{Prop}}_W(S)$, $\mathrm{Ho}^{\mathrm{Liss}}_W(S)$.

\msp{page 227}

J'ai bon espoir que le foncteur $f_!$ peut être décrit, et son adjonction avec $f^*$ établie, à aussi bon compte dans ce contexte général que dans Cat. Mais je n'ai pas envie de m'y arrêter à présent.

\section*{6 Cofibre $W$-homotopique dans $\mathrm{Cat}$.}

Par contre, j'ai envie de regarder ici sous quelles conditions sur le localiseur fondamental $W$, il existe dans $(\mathrm{Cat},W)$ une théorie duale de celle de la fibre $W$-homotopique.

La question-clef est, bien sûr, celle d'avoir une bonne prise sur les $W$-cofibrations. Or je ne connais jusqu'à présent \textit{que} les immersions ouvertes de Dwyer, les immersions fermées de Dwyer, et ce qui s'en déduit par composition etc. Le travail de \textsc{Thomason} peut motiver l'espoir que la connaissance de ces immersions-là, qui semblent ici jouer un rôle similaire (dual) à celui des morphismes $W$-parfaits, suffira pour faire une bonne théorie des `cofibres $W$-homotopiques'.

% trois traits verticaux imprimés dans la marge gauche de ce paragraphe
\marginpar{\textbar\textbar\textbar}Je suppose dans toute cette section que $W$ satisfait W(7 bis) (axiome des carrés cocartésiens), pour avoir des bonnes propriétés des immersions ouvertes de Dwyer (\footnote[124]{\begin{enumerate}
\item[a)] Ce sont des $W$-cofibrations (résulte de b) et c)).
\item[b)] Si
\end{enumerate}}), et plus généralement, pour les $W$-cofibrations. % ?? uncertain: coupure en fin de ligne "co-cartésiens" -- cocartésiens ou co-cartésiens
% foot: Version 30 mars 2003

\end{document}
